\section{Modelagem}
\label{sec:modelagem}

\subsection{Estados da autorização}

A Figura~\ref{fig:estados} representa a máquina de estados conceitual almejada.
A consulta começa em \path{QUERYING}. Produto incompatível conduz a
\path{DENIED}; ausência, ambiguidade ou código desconhecido conduz a
\path{REVIEW_REQUIRED}; e uma decisão compatível primeiro persiste a intenção
\path{PENDING_DEVICE}. Somente o reconhecimento do comando promove a
autorização para \path{AUTHORIZED}. Uma revisão excepcional permitida também
passa por essa intenção, nunca diretamente para autorização. Expiração,
conclusão, substituição ou falha retornam o dispositivo ao bloqueio.

\begin{figure}[H]
  \centering
  \noindent\resizebox{0.98\textwidth}{!}{\begin{tikzpicture}[
    node distance=8mm and 10mm,
    estado/.style={draw, rounded corners, align=center, minimum height=9mm,
      text width=25mm, fill=black!3, font=\small},
    seta/.style={-{Latex[length=1.8mm]}, thick, font=\scriptsize}
  ]
    \node[estado] (locked) {LOCKED};
    \node[estado, right=of locked] (query) {QUERYING};
    \node[estado, right=of query] (pending) {PENDING\_\\DEVICE};
    \node[estado, right=of pending] (authorized) {AUTHORIZED};
    \node[estado, right=of authorized] (completed) {COMPLETED};
    \node[estado, above=of query] (denied) {DENIED};
    \node[estado, below=of query] (review) {REVIEW\_\\REQUIRED};
    \node[estado, below=of authorized] (terminal) {EXPIRED/\\SUPERSEDED};

    \draw[seta] (locked) -- (query);
    \draw[seta] (query) -- node[above]{compatível} (pending);
    \draw[seta] (query) -- node[left]{incompatível} (denied);
    \draw[seta] (query) -- node[left]{incerto} (review);
    \draw[seta] (review) -- node[below,sloped]{override} (pending);
    \draw[seta] (pending) -- node[above]{ACK} (authorized);
    \draw[seta] (authorized) -- (completed);
    \draw[seta] (authorized) -- (terminal);
  \end{tikzpicture}}
  \caption{Máquina de estados conceitual da autorização}
  \label{fig:estados}
  \fontepropria
\end{figure}

Estados terminais, expiração, substituição e falhas conduzem ao estado físico
bloqueado; esses retornos foram omitidos da figura para preservar a
legibilidade. A implementação atual persiste \path{PENDING_DEVICE},
\path{AUTHORIZED}, \path{DENIED}, \path{REVIEW_REQUIRED},
\path{COMPLETED}, \path{EXPIRED} e \path{SUPERSEDED};
\path{QUERYING}, \path{DISPENSING}, \path{LOCKED} e \path{CANCELLED}
ainda não são estados de domínio materializados. Não existe transição de
\path{DENIED}, \path{EXPIRED}, \path{COMPLETED} ou
\path{SUPERSEDED} para \path{AUTHORIZED}. O override parte exclusivamente de
\path{REVIEW_REQUIRED} e somente quando o motivo pertence às incertezas da
consulta; falha de bomba, bico ou dispensador não pode ser contornada.

Antes de enviar \path{UNLOCK}, o backend persiste a intenção e o identificador
do comando. O gateway simulado aplica e confirma a operação de modo síncrono;
então o backend registra o evento de ACK e promove a autorização. O firmware
experimental já produz \path{COMMAND_ACK}, porém o backend ainda não consome
esse evento por MQTT. Se o bloqueio compensatório falhar, a garantia residual é
o TTL local do dispositivo, não um bloqueio imediato absoluto.

\subsection{Tabela de decisão}

\begin{table}[H]
  \centering
  \caption{Política preliminar de decisão}
  \label{tab:decisao}
  \small
  \begin{tabularx}{\textwidth}{@{}L{2.4cm}L{2.7cm}L{2.4cm}Y@{}}
    \toprule
    \textbf{Resposta} & \textbf{Produto solicitado} & \textbf{Decisão} & \textbf{Efeito físico} \\
    \midrule
    Válida e inequívoca & Compatível & Autorizar & Liberar apenas o bico solicitado por TTL. \\
    Válida e inequívoca & Incompatível & Negar & Manter todos bloqueados. \\
    Inválida ou ausente & Qualquer & Revisar & Manter bloqueado; permitir fluxo de supervisor. \\
    Código desconhecido & Qualquer & Revisar & Manter bloqueado e registrar versão da política. \\
    Válida com múltiplos & Um compatível & Autorizar & Liberar somente o produto escolhido, nunca todos. \\
    Conversão não comprovada & Qualquer & Revisar & Não inferir compatibilidade automaticamente. \\
    \bottomrule
  \end{tabularx}
  \\[2pt]\fontepropria
\end{table}

\subsection{Modelo de dados minimizado}

O modelo-alvo prevê \path{Authorization}, \path{PolicyVersion},
\path{DeviceCommand}, \path{DeviceAck}, \path{Override} e
\path{AuditEvent}. Na implementação corrente, o SQLite contém somente tabelas
de autorização e eventos de auditoria; comandos ficam no gateway em memória, e
os campos de override integram a autorização. Ainda não existem entidades
persistidas para versão da política ou confirmação do dispositivo.

A placa existe apenas em memória durante a consulta e não é chave de
relacionamento persistente. O esquema atual guarda decisão, combustível
solicitado, bomba, frentista, modo do provedor, motivo, horários e eventual
override, além de \path{client_request_id} e um fingerprint HMAC para verificar
idempotência. A consulta de compatibilidade usa cache temporário em memória,
indexado por HMAC da placa normalizada, para evitar uma segunda consulta no
mesmo fluxo. A chave HMAC fica fora do repositório; esse recurso reduz exposição,
mas é pseudonimização técnica, não anonimização. Não há colunas para placa,
Renavam, chassi, proprietário ou resposta bruta do provedor. Logs de pesquisa
serão separados dos dados operacionais e conterão somente casos sintéticos.

\pendente{validar o modelo e a enumeração oficial de códigos de combustível com
a documentação técnica liberada no ambiente autorizado.}
